\section{Experimental Design}
\label{sec:experimental-design}

Our evaluation separates measured inference components from a backend
policy replay. All GPU profiles run on one NVIDIA GeForce RTX 4090 with
24\,GB nominal device memory. The candidate-generation environment uses
PyTorch 2.11.0, Transformers 5.5.4, and CUDA 13.0, with BF16 weights.
The final checkpoint, attention backend, decoding configuration, software
versions, and source hashes are recorded in the corresponding run metadata;
they are not inferred from a model family name. Model loading and data
preparation are excluded from per-request service demands. The ASR pilot
uses a separate CTranslate2 backend and FP16 computation. GPU workloads
are scheduled sequentially so our ASR and candidate-generation profiles
do not overlap.

\subsection{Component Pilots and Development Data}

The ASR pilot uses two Chinese speech files preserved in the repository
archive, with decoded durations of 139.76 and 127.84 seconds. We decode
each file once and evaluate its first 30 seconds using cumulative
prefixes arriving at one-, two-, or four-second intervals. Three repeated
runs yield 318 timed prefix decodes after three excluded warmup calls.
Timing extends through complete consumption of the lazy ASR segment
iterator. These are offline repeated-prefix decodes, not native streaming
recognition. The recordings contain near-identical material according
to their machine transcripts, so repetitions characterize this workload
rather than adding independent scripts or speakers. We measure service
cost and transcript-prefix revision between successive updates, and
separately measure consistency with a same-model final-prefix reference;
no human transcription accuracy is claimed.

The earlier 24-utterance policy pilot is exploratory. Its outputs and
failure cases informed the research questions and cannot serve as an
unseen test of the revised system. The revised workflow instead uses
12 development cases, comprising eight positives and four negatives,
and a separately authored 40-case holdout. Development examples support
implementation and configuration choices. Before opening the holdout,
we freeze the benchmark source, prompt, model configuration, and scoring
contract, and preserve their hashes. Holdout labels and reviewer
manifests remain separate from online inputs.

\subsection{Fresh Holdout and Timing Contract}

The fresh holdout contains 32 positive utterances and eight ordinary
negative controls. Sixteen positives place their strict emphasis target
at the utterance end; sixteen retain spoken context after it. Six
positives also carry stress annotations: three contain a one-update
synthetic ASR substitution, and three require later context to resolve
an initially ambiguous interpretation. These six overlap the 32 positives.
Targets are one to four source words. Strict targets and at most two
permitted alternative spans are declared before inference. Every allowed
alternative ends at the same source-word position as its strict target,
so phrase tolerance does not move the timing reference.

Source words receive constant durations of 360, 400, or 440 milliseconds
per utterance. Partial transcripts arrive on a 500-millisecond grid.
A supplied final flag follows the first complete-text grid point by one
additional update. Revision stresses alter a target-tail word for one
update and restore it on the next, while preserving the original source
clock. These arrivals and word ends are synthetic, not measured
recognition or physical speech timing. Online methods receive only
observed partial text, arrivals, and final flags. Authored targets,
future text, safety annotations, and deadlines are used exclusively
by the evaluator.

Let $w$ denote the strict target's source-word end and $s$ its
predeclared earliest safe reveal time. The latter incorporates necessary
later context and, for a revision stress, availability of the corrected
partial. We report target coverage in both the evaluation-safe window
$[s,s+1500\,\mathrm{ms}]$ and the earlier word-end contract
$[w,w+1500\,\mathrm{ms}]$. A display before $s$ remains premature even
if it satisfies the earlier window. The two contracts are reported
explicitly rather than exchanging their reference points.

\subsection{Metrics and Measurement Boundaries}

Strict and allowed-span target coverage each use all 32 positives as
their denominator, including omissions. Negative false activation is
the number of negative clips with any displayed text divided by eight;
queued displays after the final ASR update are included. Display-level
matching precision includes negative-control displays in its denominator.
First-match latency is accompanied by its matched-case count. Paired
latency comparisons use only the common matched subset, while coverage
retains every positive case. Proposal calls, cache use, and cumulative
GPU service demand expose the computational cost of each policy.

Policy timings come from discrete-event replay with measured GPU service
demands and scripted transcript arrivals. They therefore characterize
backend scheduling under this contract. The ASR and policy pilots use
different inputs and are not combined into an end-to-end latency number.
Network delivery, browser queues, WebXR rendering, and physical headset
presentation remain outside the measured boundary.

